% !TeX root = ../../main.tex
\section{符号定义与说明}\label{cw:app:notation}

\textbf{取值空间与密度导数：}随机变量 $X,Y$ 的取值空间分别记为 $\mathcal S_X,\mathcal S_Y$（上下文讨论单一空间且无歧义时简记为 $\mathcal S$），其上可测集合族记为 $\mathcal B(\mathcal S)$。有限维实向量空间上以体积测度为默认参考测度（状态空间上的体积测度记为 $v_X$），密度导数记为 $p_X(x)=\frac{\mathrm d\mathbb P_X}{\mathrm dv_X}(x)$，条件密度简写时记为 $p(x\mid y)=p_{X\mid Y}(x\mid y)$；一般测度间的密度导数记作 $\frac{\mathrm d\mu}{\mathrm d\nu}$，绝对连续记作 $\mu\ll\nu$。数学期望与概率带分布下标写作 $\mathbb E_{X\sim\mathbb P_X}[\cdot]$ 与 $\mathbb P_X(\cdot)$。

\textbf{概率单纯形与全一向量：}对正整数 $n$，$\mathbb R_+=[0,\infty)$，定义
\begin{equation*}
\Delta^{n-1}:=\left\{\bm w\in\mathbb R_+^n:\sum_{i=1}^n w_i=1\right\},\qquad
\mathbf1_n:=(1,\dots,1)^{\mathrm T}.
\end{equation*}
$\Delta^{n-1}$ 是 $n$ 个非负概率权重组成的集合，上标表示其维数；允许零权重，严格正权重另行说明。

\textbf{单点质量测度与示性函数：}对 $x\in\mathcal S$ 和可测集合 $A\in\mathcal B(\mathcal S)$，
\begin{equation*}
\delta_x(A):=\mathbb I_A(x)=
\begin{cases}1,&x\in A,\\0,&x\notin A,\end{cases}
\qquad \int_{\mathcal S}f(s)\,\mathrm d\delta_x(s)=f(x).
\end{equation*}
$\delta_x$ 是集中于单点 $x$ 的概率测度，与表示概率单纯形的大写 $\Delta$ 不同。

\textbf{概率测度与经验测度：}$\mathcal M(\mathcal S)$ 表示 $\mathcal S$ 上的全体非负测度，概率测度全体记为 $\mathcal M^1(\mathcal S):=\{\mu\in\mathcal M(\mathcal S):\mu(\mathcal S)=1\}$，不得将两者混用。记单点质量测度行数组 $\delta_{x_{1..N}}:=(\delta_{x_1},\dots,\delta_{x_N})$，对权重列向量 $\bm w\in\Delta^{N-1}$，加权与等权经验测度分别记为
\begin{equation*}
\delta_{x_{1..N}}\bm w:=\sum_{i=1}^Nw_i\delta_{x_i},\qquad
\hat{\mathbb P}_{N,\bm w;X}:=\delta_{x_{1..N}}\bm w,\qquad
\hat{\mathbb P}_{N;X}:=\frac1N\delta_{x_{1..N}}\mathbf1_N=\frac1N\sum_{i=1}^N\delta_{x_i}.
\end{equation*}
这里 $x_i$ 是样本实现值，行数组右乘列向量表示测度的有限凸组合。

\textbf{索引排列与滤波阶段：}用 $..$ 表示行排列数组 $x_{1..N}:=(x_1,\dots,x_N)$，用 $:$ 表示列拼接序列 $y_{1:k}:=(y_1,\dots,y_k)^{\mathrm T}$。对样本随机列向量 $X_i\in\mathbb R^{d\times1}$，记
\begin{equation*}
X_{1..N}:=(X_1,\dots,X_N)\in\mathbb R^{d\times N},\qquad
\mathbb P_{X_{1..N}}\text{ 为该样本族的联合分布}.
\end{equation*}
一般二元随机变量的联合分布记为 $\mathbb P_{(X,Y)}$。
上标 $-$ 与 $+$ 分别表示预测与后验阶段。位置变换中的 $X_k$ 表示按列拼接的粒子实现矩阵，每列是一个 $n_x$ 维粒子：
\begin{equation*}
\begin{aligned}
&X_k:=x_k^{(1..N)}:=[x_k^{(1)},\dots,x_k^{(N)}]\in\mathbb R^{n_x\times N},\\
&X_k^\pm:=x_k^{\pm,(1..N)}\in\mathbb R^{n_x\times N}.
\end{aligned}
\end{equation*}
$(\bm T_k)_{ij}$ 表示源粒子 $i$ 对目标粒子 $j$ 的贡献系数，故 $X_k^+=X_k^-\bm T_k$，$\bm T_k=N\bm\Pi_k$。概率分布下标中的 $X_k$ 仍表示状态随机变量，由上下文区分。

\textbf{推前、拉回与对偶配对：}对可测映射 $T:\mathcal S_X\to\mathcal S_Y$，
\begin{equation*}
\begin{aligned}
&(T_*\mu)(B):=\mu(T^{-1}(B)),\qquad T^*f:=f\circ T,\\
&\langle f,\mu\rangle_{\mathcal S}:=\int_{\mathcal S}f(s)\,\mathrm d\mu(s),\qquad
\langle f,T_*\mu\rangle_{\mathcal S_Y}=\langle T^*f,\mu\rangle_{\mathcal S_X}.
\end{aligned}
\end{equation*}
其中 $B\in\mathcal B(\mathcal S_Y)$，函数与测度的对偶配对尖括号须标明测度所在空间，积分对非负可测或相应可积函数使用。$\pi_i$ 表示取第 $i$ 个分量的坐标投影，$\mathbb P_X=X_*\mathbb P$。

\textbf{核复合与测度-核乘积耦合：}对条件转移核 $\kappa:\mathcal S_X\to\mathcal M(\mathcal S_Y)$ 与测度 $\mu\in\mathcal M(\mathcal S_X)$，复合测度记作 $\kappa\circ_m\mu:=\int_{\mathcal S_X}\kappa(x)\,\mu(\mathrm dx)\in\mathcal M(\mathcal S_Y)$；测度-核乘积耦合记作 $\mu\otimes_m\kappa\in\mathcal M(\mathcal S_X\times\mathcal S_Y)$，由此联合分布写作 $\mathbb P_{(X,Y)}=\mathbb P_X\otimes_m\mathbb P_{Y\mid X}$，独立乘积测度 $\mu\otimes\nu$（如 $\mathbb P_X\otimes\mathbb P_Y$）为其特例。

\textbf{联合概率矩阵与条件概率矩阵：}在正质量离散支撑上，记
\begin{equation*}
\begin{aligned}
&a_i:=\mathbb P_X(\{x_i\}),\qquad b_j:=\mathbb P_Y(\{y_j\}),\qquad
\Pi_{ij}:=\mathbb P_{(X,Y)}(\{(x_i,y_j)\}),\\
&P_{Y\mid X;i,j}:=\mathbb P_{Y\mid X=x_i}(\{y_j\})=\Pi_{ij}/a_i,\qquad
P_{X\mid Y;j,i}:=\mathbb P_{X\mid Y=y_j}(\{x_i\})=\Pi_{ij}/b_j.
\end{aligned}
\end{equation*}
条件矩阵的行对应条件变量，列对应被条件化变量；$\bm\Pi=(\Pi_{ij})$ 是联合概率矩阵。$\Pi(\mathbb P_X,\mathbb P_Y)$ 表示具有指定边缘的耦合测度集合，$U(\bm a,\bm b)$ 为其离散矩阵表示，见式~\eqref{cw:eq:coupling-family}、\eqref{cw:eq:transport-polytope}。

\textbf{矩阵与向量运算：}矩阵转置写作 $\bm A^{\mathrm T}$，迹写作 $\operatorname{tr}(\bm A)$，$\operatorname{diag}(\bm u)$ 将向量 $\bm u$ 置于矩阵对角线上；$\odot$、$\oslash$ 分别表示逐元素乘、除，$\otimes_k$ 表示矩阵分块直积。实矩阵迹内积记为 $\langle\bm A\mid\bm B\rangle:=\operatorname{tr}(\bm A^{\mathrm T}\bm B)=\sum_{i,j}A_{ij}B_{ij}$；向量一范数记为 $\|\bm z\|_1:=\sum_i|z_i|$，$\|\bm z\|$ 表示向量二范数。

\textbf{输运代价与熵正则化：}离散代价矩阵记为 $\bm C=(C_{ij})$，其中 $C_{ij}:=c(x_i,y_j)$。相对熵记为 $D_{\mathrm{KL}}(\mu\|\nu):=\int\log(\mathrm d\mu/\mathrm d\nu)\,\mathrm d\mu$（$\mu\ll\nu$ 时），否则为 $+\infty$；$\varepsilon>0$ 为熵正则化参数，$\tau>0$ 为边缘残差容差。

\textbf{平方可积空间与正交投影：}取值于 $\mathbb R^{n_x}$ 的平方可积随机变量空间记为 $\mathcal H:=L^2(\Omega,\mathcal F,\mathbb P;\mathbb R^{n_x})$；闭线性子空间关系记为 $\mathcal W\leqslant\mathcal H$，向 $\mathcal W$ 的正交投影算子与余投影算子分别记为 $\mathscr P_{\mathcal W}$ 与 $\mathscr P_{\mathcal W}^{\perp}$。

\textbf{参数空间、对偶伴随与压缩还原算子：}模型参数随机变量 $\bm\Theta$ 的取值空间记为 $\mathcal S_{\bm\Theta}$，参数实现记为 $\bm\theta\in\mathcal S_{\bm\Theta}$，含参条件密度记为 $p_{Y\mid(Z,\bm\Theta)}(y\mid z,\bm\theta)$。赋范或内积空间上的连续线性对偶与内积伴随算子统一记作上标 ${}^\star$，优化问题的最优解统一记作空心星上标 ${}^{\optstar}$（如 $\bm\theta^{\optstar},\mathscr T_{\Rightarrow}^{\optstar}$），以严格区分伴随运算与最优性。高维观测空间 $\mathcal S_X=\mathbb R^{n_D}$ 与低维潜空间 $\mathcal S_Z=\mathbb R^{n_d}$（$n_d<n_D$）之间的压缩算子与还原算子分别记为 $\mathscr T_{\Rightarrow}:\mathcal S_X\to\mathcal S_Z$ 与 $\mathscr T_{\Leftarrow}:\mathcal S_Z\to\mathcal S_X$，线性情形下二者在标准基下的表示矩阵记为 $\bm T_{\Rightarrow},\bm T_{\Leftarrow}$。

